\title{xLSTM: Extended Long Short-Term Memory}

\begin{abstract}
The foundational concepts behind Long Short-Term Memory (LSTM) networks—namely, the constant error carousel and gating mechanisms—emerged in the 1990s. LSTMs have since proven to be enduring and instrumental within deep learning, playing a key role in the development of the initial Large Language Models (LLMs). Nevertheless, the introduction of the Transformer model, defined by its parallelizable self-attention, signified a major breakthrough, propelling model performance well beyond that of LSTMs, especially at large scales. This context motivates a fundamental inquiry: To what extent can LSTMs advance in language modeling if scaled up to billions of parameters, while harnessing state-of-the-art methods from contemporary LLMs and addressing well-known LSTM shortcomings? We answer this by first proposing exponential gating, complemented by suitable normalization and stabilization strategies. Next, we redesign LSTM memory, leading to: (i) sLSTM, which features scalar memory, scalar updates, and an innovative mixing process; (ii) mLSTM, a fully parallelizable variant utilizing matrix-based memory along with a covariance-driven update mechanism. Embedding these advanced LSTM designs into residual block frameworks results in xLSTM blocks, which are subsequently layered residually to build xLSTM architectures. Through the incorporation of exponential gating and refined memory approaches, xLSTM models demonstrate improved scaling properties and competitive performance relative to leading-edge Transformers and State Space Models.\\
Code available at: \url{https://github.com/NX-AI/xlstm}
\end{abstract}

\section{Introduction}
\label{sec:intro}

The Long Short-Term Memory (LSTM) ideas~\citep{Hochreiter:91,Hochreiter:97nips,Hochreiter:97}, i.e., the constant error carousel and gating, were introduced to overcome the vanishing gradient problem of recurrent neural networks~\citep{Hochreiter:91,Hochreiter:00book}:

\begin{align}
\label{eq:lstm_idea}
  c_t \ = \ f_t \ c_{t-1} \ + \ 
          i_t \  z_t \ , \quad  
          h_t \ = \ o_t \ \psi({c_t}) \ . 
\end{align}

The constant error carousel functions by incrementally adjusting the previous cell state $c_{t-1}$ (shown in green) through the integration of cell inputs $z_t$, with this process regulated via sigmoid gates (depicted in blue). The input gate $i_t$ and forget gate $f_t$ manage these updates, whereas the output gate $o_t$ governs the memory cell’s output—specifically, the hidden state $h_t$. After normalization or compression by $\psi$, the cell state passes through the output gate to yield the hidden state.

LSTMs have achieved widespread success across multiple fields \citep{Hochreiter:01,Hochreiter:07,Schmidhuber:15}, dominating text generation tasks until the emergence of Transformer architectures in 2017~\citep{Vaswani:17}. Their proficiency has been established in numerous sequential tasks, including but not limited to text generation~\citep{Graves:13,Karpathy:15blog}, handwriting generation~\citep{Graves:13}, sequence-to-sequence translation~\citep{Sutskever:14nips}, program evaluation~\citep{Zaremba:14arxiv}, image captioning~\citep{Karpathy:15,Hossain:19}, code generation~\citep{Karpathy:15blog}, rainfall-runoff simulation~\citep{Kratzert:18,Kratzert:19}, and hydrological prediction for flood warnings~\citep{Nearing:24}. Within the field of reinforcement learning, LSTM-based models represent the top-performing sequence architectures, as shown by applications such as AlphaStar for StarCraft II~\citep{Vinyals:17short}, OpenAI Five for Dota 2~\citep{OpenAI:19}, and controllers used in nuclear fusion devices~\citep{Degrave:22short}. One of their distinguishing features is their ability to learn and retain abstract representations by efficiently extracting semantic content into their memory components~\citep{Karpathy:15blog}, which is exemplified in the identification of number and syntax neurons~\citep{Lakretz:19}, linguistic processing units~\citep{Bau:19}, and sentiment tracking neurons~\citep{Radford:17arxiv}. Despite advances in other architectures, LSTMs continue to be employed in high-impact scenarios~\citep{Degrave:22short,Nearing:24} and have demonstrated enduring robustness.

While LSTMs have achieved notable milestones, they present three critical shortcomings: (i) Inability to correct previous storage selections. This can be highlighted with the 
{\it Nearest Neighbor Search} scenario (further illustrated in Appendix~\ref{sec:appNearestNeighborSearch}): when given a reference vector, the task involves sequentially traversing a sequence to identify the most similar vector and retrieving its associated value by the sequence's end. The left panel in Figure~\ref{fig:lstmProblems} displays mean squared error for this problem. Standard LSTMs have difficulty updating a previously stored item if a more similar vector appears later in the sequence, an issue mitigated by our xLSTM, which employs exponential gating. (ii) Restricted capacity for information retention due to compression of information into scalar cell states. This limitation becomes apparent in the context of {\it Rare Token Prediction}. According to the right panel of Figure~\ref{fig:lstmProblems}, which reports perplexity for token prediction on Wikitext-103~\citep{Merity:17} across different token frequency bands, LSTM underperforms on low-frequency tokens—largely because of its storage constraints. Our xLSTM addresses this with its adoption of matrix memory. (iii) Poor parallelization arising from memory interdependence—that is, the recurrent hidden-to-hidden connections from one timestep to the next necessitate sequential computation. 

Collectively, such drawbacks of LSTM have contributed to the ascendancy of Transformers~\citep{Vaswani:17} within language modeling. The question therefore arises: What levels of performance in language modeling become achievable if we eliminate these drawbacks and scale LSTMs to the parameter magnitudes of current Large Language Models?

\section{Extended Long Short-Term Memory}
\label{sec:xLSTM}

In order to address the shortcomings of standard LSTMs, Extended Long Short-Term Memory (xLSTM) incorporates two principal innovations into the conceptual framework of Equation~\eqref{eq:lstm_idea}: exponential gating mechanisms and distinctive memory architectures. These enhancements introduce two distinct variants to the LSTM family: (i) the sLSTM, detailed in Section~\ref{sec:sLSTM}, which utilizes a scalar memory, scalar updates, and allows for memory mixing, and (ii) the mLSTM, as described in Section~\ref{sec:mLSTM}, featuring a matrix-based memory system and an update mechanism based on covariance (outer products), enabling complete parallelization. Exponential gating is utilized in both sLSTM and mLSTM to augment the functionalities of conventional LSTMs. To facilitate parallel processing, mLSTM forgoes memory mixing, that is, recurrent connections between hidden states. Both variants can accommodate multiple memory cells; for sLSTM, this involves mixing memory across cells, while mLSTM maintains equivalence between multiple heads and cells. Additionally, the sLSTM architecture can include multiple heads, where mixing occurs exclusively among cells within individual heads, not between heads, leading to a novel approach to memory mixing facilitated by both multiple heads and exponential gating. 

Embedding these alternative LSTM mechanisms into the residual block framework gives rise to xLSTM blocks (refer to Section~\ref{sec:xLSTMblock}). Constructing network architectures through residual stacking of xLSTM blocks yields complete xLSTM architectures (see Section~\ref{sec:xLSTMarchitecure}). Figure~\ref{fig:xlstm_sketch} presents a diagram outlining the xLSTM architecture and its constituent elements.

\subsection{Review of the Long Short-Term Memory}
\label{sec:orgLSTM}

The foundational concept of LSTM~\citep{Hochreiter:91,Hochreiter:97nips,Hochreiter:97} centers around a scalar memory cell, which serves both as a computational and storage mechanism, thereby mitigating the issue of vanishing gradients~\citep{Hochreiter:91,Hochreiter:00book} via the constant error carousel realized by the cell state update. The architecture of this memory cell encompasses three gates: input, output, and forget. The forget gate was later incorporated by \citet{Gers:00}. 
The following expressions define the LSTM memory cell's update procedures at the $t$-th time step:

\begin{align}
c_t \ &= \  f_t \ c_{t-1} \ + \ i_t \ z_t &  & &\text{cell state} \\
h_t  \ &= \ o_t \ \tilde{h}_t \ , & \tilde{h}_t \ &= \ \psi \left( c_t\right)
  &\text{hidden state} \\
z_t \ &= \ \varphi \left( \tilde{z}_t \right) \ , 
  &\tilde{z}_t \ &=  \ \mathbf{w}^\top_{z} \ \mathbf{x}_t \ + \
  r_{z}  h_{t-1} \ + \  b_{z} \ \
  &\text{cell input} \\
i_t \ &= \ \sigma \left( \tilde{i}_t  \right) \ , 
  &\tilde{i}_t \ &= \ \mathbf{w}^\top_{i} \ \mathbf{x}_t \ + \
  r_{i}  \ h_{t-1} \ + \  b_{i} \ \
  &\text{input gate} \\
f_t \ &= \ \sigma \left(  \tilde{f}_t \right) \ , 
  &\tilde{f}_t \ &= \ \mathbf{w}^\top_{f} \ \mathbf{x}_t  \ + \
  r_{f}  \ h_{t-1} \ + \  b_{f} \ \
  &\text{forget gate} \\
o_t \ &= \ \sigma \left( \tilde{o}_t \right) \ , 
  &\tilde{o}_t  \ &= \ \mathbf{w}^\top_{o} \ \mathbf{x}_t \ + \
  r_{o}  \ h_{t-1} \ + \  b_{o} \ \
  &\text{output gate} 
\end{align}

The vectors $\mathbf{w}_{z}$, $\mathbf{w}_{i}$, $\mathbf{w}_{f}$, and $\mathbf{w}_{o}$ represent the input weights that link the input $\mathbf{x}_t$ to the cell input, input gate, forget gate, and output gate, respectively. Likewise, the terms $r_{z}$, $r_{i}$, $r_{f}$, and $r_{o}$ denote recurrent weights connecting the previous hidden state $h_{t-1}$ to the corresponding components: cell input, input gate, forget gate, and output gate. The biases for each component are denoted as $b_{z}$, $b_{i}$, $b_{f}$, and $b_{o}$. The functions $\varphi$ and $\psi$ are applied as activation functions to the cell input and the hidden state, with $\tanh$ typically used in both cases. To prevent the cell state from growing without bounds, $\psi$ is employed for normalization or squashing purposes. Sigmoid functions serve as the activation mechanisms for all gates, defined as $\sigma \left( x \right) = 1/(1+\exp(-x))$. In more advanced architectures, memory was represented as a vector of multiple scalar cells $\mathbf{c}_t \in \mathbb{R}^{d}$, enabling the application of recurrent weight matrices $\mathbf{R} \in \mathbb{R}^{d \times d}$ that interconnect outputs across the set of memory cells~\citep{Greff:15}; further technical details can be found in Appendix~\ref{sec:apporgLSTM}. Results from ablation studies indicate that every part of the memory cell plays an essential role~\citep{Greff:15}.

\subsection{sLSTM}
\label{sec:sLSTM}

To enhance LSTMs with the capability to alter their storage choices, we implement exponential gates (depicted in red), along with mechanisms for normalization and stabilization. Specifically, the activation functions of both input and forget gates can be exponential. For the purpose of normalization, we propose a normalizer state that accumulates the sum of the input gate multiplied by each subsequent forget gate.

The scalar sLSTM forward pass is:

\begin{align}
\label{eq:slstmforward}
c_t \ &= \  f_t \ c_{t-1} \ + \ i_t \ z_t & & 
 &\text{cell state} \\
n_t \ &= \  f_t \ n_{t-1} \ + \ i_t  & & 
 &\text{normalizer state} \\
h_t \ &= \ o_t \ \tilde{h}_t \ ,
  &\tilde{h}_t \ &= \ c_t / n_t
&\text{hidden state} \\
z_t \ &= \ \varphi \left( \tilde{z}_t \right) \ , 
  &\tilde{z}_t \ &=  \ \mathbf{w}^\top_{z} \ \mathbf{x}_t \ + \
  r_{z}  h_{t-1} \ + \  b_{z}
  &\text{cell input} \\
i_t \ &= \ \!\exp \left( \tilde{i}_t  \right) \ , 
  &\tilde{i}_t \ &= \ \mathbf{w}^\top_{i} \ \mathbf{x}_t \ + \
  r_{i}  \ h_{t-1} \ + \  b_{i}
  &\text{input gate} \\
f_t \ &= \ \sigma \left(  \tilde{f}_t \right) \ \text{OR} \
\!\exp \left(  \tilde{f}_t \right) \ ,
  &\tilde{f}_t \ &= \ \mathbf{w}^\top_{f} \ \mathbf{x}_t  \ + \
  r_{f}  \ h_{t-1} \ + \  b_{f}
  &\text{forget gate} \\
o_t \ &= \ \sigma \left( \tilde{o}_t \right) \ , 
  &\tilde{o}_t  \ &= \ \mathbf{w}^\top_{o} \ \mathbf{x}_t \ + \
  r_{o}  \ h_{t-1} \ + \  b_{o}
  &\text{output gate} 
\end{align}

We transfer the original LSTM gating techniques, i.e., input- and/or hidden-dependent gating plus bias term, to the new architectures. Exponential activation functions can lead to large values that cause overflows. Therefore, we stabilize gates with an additional state \mbox{$m_t$~\citep{Milakov:18arxiv}}:

\begin{align}
\label{eq:slstmstabil}
m_t \ &= \ \max \left( \log ( f_t ) + m_{t-1} , \log ( i_t ) \right) &\text{stabilizer state} \\
i'_t \ &= \ \!\exp \left( \log \left ( i_t \right) - m_t \right) \ = \ \!\exp \left( \tilde{i}_t - m_t \right) \ \, 
  &\text{stabil. input gate} \\
f'_t \ &= \ \!\exp \left( \log \left( f_t \right) + m_{t-1} - m_t \right) \ \, 
  &\text{stabil. forget gate}
\end{align}

Appendix~\ref{sec:appsLSTM} demonstrates that substituting $f_t$ with $f'_t$ and $i_t$ with $i'_t$ during the forward computation preserves both the network's overall output and the gradients of the loss function relative to the parameters.

\paragraph{New Memory Mixing.} Similar to the classical LSTM architecture, sLSTM is capable of supporting several memory cells concurrently (refer to Appendix~\ref{sec:appsLSTM}). The presence of multiple memory cells facilitates the process of memory mixing through recurrent links $\mathbf{R}_{\mathbf{z}}$, $\mathbf{R}_{\mathbf{i}}$, $\mathbf{R}_{\mathbf{f}}$, $\mathbf{R}_{\mathbf{o}}$ originating from the hidden state vector $\mathbf{h}$ to the memory cell input $\mathbf{z}$ and to each of the gating variables $\mathbf{i}$, $\mathbf{f}$, and $\mathbf{o}$. A novel feature in this form of memory mixing is attributed to exponential gating. In sLSTM, multiple heads can exist, allowing mixing of memories internally within a head, while preventing interactions between different heads. By combining head-based organization in sLSTM with exponential gating mechanisms, a fundamentally new strategy for memory mixing is introduced.

\subsection{mLSTM}
\label{sec:mLSTM}

In order to improve the memory capability of LSTMs, we substitute the traditional scalar memory cell $c \in \mathbb{R}$ with a matrix-form cell $\mathbf{C} \in \mathbb{R}^{d \times d}$. This modification enables retrieval operations to be conducted using matrix multiplication. At a given time step $t$, the model stores two vectors: a key vector $\mathbf{k}_t \in \mathbb{R}^d$ and a value vector $\mathbf{v}_t \in \mathbb{R}^d$ (adopting terminology from the Transformer architecture). Subsequently, at time $t + \tau$, the value $\mathbf{v}_t$ is accessed using a query vector $\mathbf{q}_{t+\tau} \in \mathbb{R}^d$. This configuration aligns with the framework of Bidirectional Associative Memories (BAMs) \citep{Kohonen:72,Anderson:72,Nakano:72,Anderson:77}.

The covariance update rule~\citep{Sejnowski:77,Dayan:91} for storing a key-value pair is

\begin{align}
\mathbf{C}_t \ &= \ \mathbf{C}_{t-1} \ + \ \mathbf{v}_{t} \ \mathbf{k}_t^\top \ .
\end{align}

We posit the application of layer normalization prior to mapping the input data to keys and values, ensuring these vectors have zero mean. The covariance update rule achieves optimality~\citep{Dayan:91} in maximizing the separability of retrieved binary codes, which aligns with maximizing the signal-to-noise ratio. When restricting memory retrieval to pairwise operations and accepting quadratic computational complexity as in attention mechanisms~\citep{Krotov:16,Krotov:17,Ramsauer:21}, even greater separability can be attained. Notably, the covariance update rule shares its formulation with Fast Weight Programmers~\citep{Schmidhuber:92ncfastweights,Schlag:21}, for which later adaptations incorporate a constant decay factor (applied to $\mathbf{C}_{t-1}$) as well as a constant learning rate (applied to 
$\mathbf{v}_{t} \mathbf{k}_t ^\top$)~\citep{Ba:16}. Motivated by this, we embed the covariance update procedure within the LSTM architecture: the forget gate modulates the decay rate, the input gate controls the learning rate, and the output gate scales the returned memory vector.

For this matrix memory, the normalizer state is the weighted sum of key vectors, where each key vector is weighted by the input gate and all future forget gates. Again, the normalizer state keeps record of the strength of the gates. Since the dot product between query and normalizer state can be close to zero, we use the absolute value of this dot product and lower bound it by a threshold (typically 1.0) as done previously~\citep{Sun:23arxiv}. The mLSTM forward pass is:

\begin{align}
\label{eq:mlstm_recurrent_begin}
\mathbf{C}_t \ &= \  f_t \ \mathbf{C}_{t-1} \ + \ 
  i_t \ \mathbf{v}_t \ \mathbf{k}_t^\top & &  &\text{cell state} \\
\mathbf{n}_t \ &= \  f_t \ \mathbf{n}_{t-1} \ + \ 
  i_t \ \mathbf{k}_t & &  &\text{normalizer state} \\
\mathbf{h}_t  \ &= \ \mathbf{o}_t \ \odot \ \tilde{\mathbf{h}}_t
\ , \qquad \qquad 
\tilde{\mathbf{h}}_t \ = \   \mathbf{C}_t \mathbf{q}_t \ / \ 
  \max \left\{ |\mathbf{n}_t^\top \mathbf{q}_t|, 1 \right\} 
   &&&\text{hidden state} \\
\mathbf{q}_t \ &= \ \mathbf{W}_q \ \mathbf{x}_t \ + \ \mathbf{b}_q  & &  &\text{query input} \\
\mathbf{k}_t \ &= \ \frac{1}{\sqrt{d}} \mathbf{W}_k \ \mathbf{x}_t \ + \ \mathbf{b}_k  & &  &\text{key input} \\
\mathbf{v}_t \ &= \ \mathbf{W}_v \ \mathbf{x}_t \ + \ \mathbf{b}_v  & &  &\text{value input} \\
i_t \ &= \ \!\exp \!  \! \left( \tilde{i}_t  \right) \ , 
 \qquad \qquad \ \ \ \ \,
  \tilde{i}_t \ = \ \mathbf{w}^\top_{i} \ \mathbf{x}_t \ + \  b_{i} \ \ \ 
  &&&\text{input gate} \\
f_t \ &= \ \sigma \!  \left(  \tilde{f}_t \right) \ \text{OR} \ \!\exp \!  \! \left(  \tilde{f}_t \right) , \ \ \: \!
\tilde{f}_t \ = \ \mathbf{w}^\top_{f} \ \mathbf{x}_t  \ + \
  b_{f}
  &&& \text{forget gate} \\
\label{eq:mlstm_recurrent_end}
\mathbf{o}_t \ &= \ \sigma \left( \tilde{\mathbf{o}}_t \right) \ , \qquad \qquad \quad \ \ \ \,
  \tilde{\mathbf{o}}_t  \ = \ \mathbf{W}_{\mathbf{o}} \ \mathbf{x}_t \ + \
  \mathbf{b}_{\mathbf{o}}
  &&&\text{output gate} 
\end{align}

The mLSTM architecture allows for several memory cells, similar to the standard LSTM design. With mLSTM, utilizing multiple heads is functionally the same as employing multiple cells, due to the absence of memory interaction between them. To mitigate instability arising from the exponential gating mechanisms in mLSTM, we adopt the stabilization approach used in sLSTM, outlined in Equation~\ref{eq:slstmstabil}. 
Because mLSTM operates without memory mixing, its recurrence relation can be re-expressed in a parallelized format. For further information, please consult Appendix~\ref{sec:appmLSTM}.

\subsection{xLSTM Architecture}
\label{sec:xLSTMarch}

\paragraph{xLSTM Blocks.}
\label{sec:xLSTMblock}
An xLSTM block aims to produce a non-linear, high-dimensional summary of preceding elements to enhance the differentiation between varying historical contexts. This clear distinction among past histories is essential for accurate sequence element prediction, such as forecasting the subsequent token. To achieve this, we leverage Cover’s Theorem~\citep{Cover:65}, which asserts that non-linear embeddings into higher-dimensional spaces improve the odds of linear separation compared to the original representation space. We analyze two residual block structures: (i) A residual block incorporating a post up-projection mechanism, analogous to the Transformer approach, wherein the sequence history is non-linearly condensed within the original dimensionality, then projected linearly into a higher-dimensional space, followed by a non-linear activation, and subsequently mapped back to the initial space via a linear transformation; this design is illustrated in the left panel of Figure~\ref{fig:Backbones}
and is represented in the third column of Figure~\ref{fig:xlstm_sketch}. 
A comprehensive depiction is available in Figure~\ref{fig:appsLSTM_detailed}
in the appendix. (ii) A residual block that applies a pre up-projection, similar to architectures in State Space Models, initially performing a linear up-projection into a higher-dimensional space,

captures previous information in a non-linear fashion within a high-dimensional space, and subsequently applies a linear transformation to revert to the original space. If an xLSTM block is built with an sLSTM, a post up-projection block is applied. Conversely, when utilizing an mLSTM within the xLSTM, a pre up-projection block is preferred, owing to the expanded memory available in the higher-dimensional representation. For further clarification, refer to the left panel of Figure~\ref{fig:Backbones}, the third column of Figure~\ref{fig:xlstm_sketch}, or Figure~\ref{fig:appsLSTM_detailed} in the appendix.

\paragraph{xLSTM Architecture. }
\label{sec:xLSTMarchitecure}
The xLSTM model is formed by stacking component blocks using residual connections, following the approach of~\citep{Srivastava:15,He:16}. We adopt pre-LayerNorm residual backbones~\citep{Ba:16b}, consistent with the architecture of modern Large Language Models. This structure is visualized in the final column of Figure~\ref{fig:xlstm_sketch}.

\subsection{Memory and Speed Considerations}

In contrast to Transformers, xLSTM architectures maintain both linear computational requirements and fixed memory usage irrespective of the input sequence length. Owing to its compressive memory design, xLSTM is particularly advantageous for industrial use cases and deployments on edge devices.

The memory utilized by mLSTM, while parameter-free, necessitates intensive computation because of its $d \times d$ matrix storage and associated $d \times d$ updates. Here, computational burden is increased in exchange for greater memory capacity. Despite this, leveraging the parallel processing capabilities of GPUs mitigates the impact on actual runtime, causing only negligible differences in wall clock performance.

mLSTM models offer parallelization abilities similar to FlashAttention~\citep{Dao:22,Dao:23} and GLA~\citep{Yang:23arxiv}, unlike sLSTM, where memory mixing via hidden-to-hidden connections precludes such parallel computation. To address this, we have engineered a highly efficient CUDA-based implementation for sLSTM, incorporating optimizations at the GPU register level, which ensures that sLSTM achieves speeds within roughly a factor of two compared to mLSTM.

\section{Related Work}
\label{sec:related}

\paragraph{Linear Attention.} A number of approaches have been proposed to address the Transformer's quadratic complexity with respect to input sequence length, aiming to achieve linear complexity. Synthesizer eliminates token-to-token dependencies by learning synthetic attention distributions~\citep{Tay:20arxiv}. Linformer introduces self-attention via low-rank factorization, providing a linear approximation to the original mechanism~\citep{Wang:20arxiv}. Linear Transformer redefines the attention process to be inherently linear~\citep{Katharopoulos:20}. Performer makes use of positive orthogonal random features to enable a linear-time softmax attention approximation~\citep{Choromanski:21}. Furthermore, fast long convolution operations are used instead of traditional attention in Structured Global Convolution (SGConv)~\citep{Li:22} and the Hyena Hierarchy~\citep{Poli:23}.

\paragraph{State Space Models.} In recent years, State Space Models (SSMs) have surged in popularity due to their linear complexity relative to sequence length and their competitive results when compared to Transformers. The introduction of the Structured State Space sequence model (S4)~\citep{Gu:21} marked a foundational moment, subsequently inspiring developments such as the Diagonal State Space (DSS) model~\citep{Gupta:22}, Gated State Space (GSS) models~\citep{Mehta:22}, S5 model~\citep{Smith:22}, Bidirectional Gated SSM (BiGS)~\citep{Wang:22}, the H3 model~\citep{Fu:23}, and Mamba~\citep{Gu:24arxiv}.

\paragraph{Recurrent Neural Networks.} Within the domain of large-scale language modeling, Recurrent Neural Networks (RNNs) have emerged as an alternative to Transformer-based architectures and attention mechanisms, primarily owing to their linear scaling with sequence length. Models utilizing Deep Linear Recurrent Units (LRUs) have demonstrated encouraging outcomes in language modeling tasks~\citep{Orvieto:23, De:24arxiv}, while variants such as the Hierarchically Gated Linear RNN (HGRN)~\citep{Qin:23} and its successor HGRN2~\citep{Qin:24arxiv} have similarly exhibited favorable results. Among RNN approaches, RWKV~\citep{Peng:23arxivshort,Peng:24arxivshort} is particularly notable, achieving performance levels comparable to those of Transformer models.

\paragraph{Gating.}
Gating is a fundamental mechanism introduced by LSTM and has since been adapted and reconceptualized in numerous contemporary models. This technique plays a central role in architectures such as HGRN~\citep{Qin:23}, HGRN2~\citep{Qin:24arxiv}, Gated Linear Attention (GLA)~\citep{Yang:23arxiv}, Gated State Space (GSS) models~\citep{Mehta:22}, Bidirectional Gated SSM (BiGS)~\citep{Wang:22}, Moving Average Equipped Gated Attention (MEGA)~\citep{Ma:22}, RWKV~\citep{Peng:23arxivshort}, and Mamba~\citep{Gu:24arxiv}.

\paragraph{Covariance Update Rule.} To improve the storage properties of the mLSTM unit, we integrated a matrix-based memory featuring a covariance update procedure. Comparable update strategies are incorporated in Fast Weight Programmers \citep{Schmidhuber:92ncfastweights,Schlag:21}, RWKV-5 and RWKV-6~\citep{Peng:24arxivshort}, Retention~\citep{Sun:23arxiv}, Linear Transformer~\citep{Katharopoulos:20}, and HGRN2~\citep{Qin:24arxiv}.

\paragraph{Most Related.} Among existing architectures, Retention~\citep{Sun:23arxiv}, RWKV~\citep{Peng:23arxivshort,Peng:24arxivshort}, and HGRN2~\citep{Qin:24arxiv} are the most conceptually similar to xLSTM, as they all incorporate the mechanisms of matrix memory and gating. Yet, these models differ from the novel sLSTM in a fundamental way: they lack support for memory mixing. The ability to mix memory is crucial for addressing state tracking challenges, rendering LSTMs more powerful in terms of expressivity than State Space Models (SSMs) and Transformers~\citep{Merrill:24,Deletang:23}. Tasks such as code execution or entity tracking in lengthy texts depend on effective state tracking.

\paragraph{Residually Stacking Architectures.} Similar to most state-of-the-art deep neural network models, the xLSTM framework is designed by sequentially stacking its component modules with residual connections~\citep{Srivastava:15,He:16}.

This methodology was instrumental in facilitating the development of deep convolutional architectures~\citep{He:16} and Transformer models~\citep{Vaswani:17}. Transformers serve as the foundational technology powering Large Language Models (LLMs) such as GPT-3~\citep{Brown:20short}, ChatGPT~\citep{Schulman:22short}, GPT-4~\citep{Achiam:23gpt4short}, Megatron-LM~\citep{Shoeybi:19}, Gopher~\citep{Rae:21short}, ERNIE 3.0 Titan~\citep{Wang:21arxivshort}, GLaM~\citep{Du:21glamshort}, Chinese M6~\citep{Lin:21}, AlexaTM 20B (multilingual)~\citep{Soltan:22}, OPT~\citep{Zhang:22opt}, Chinchilla~\citep{Hoffmann:22short}, BLOOM~\citep{Scao:22arxivshort}, GLM-130B~\citep{Zeng:22arxivshort}, LaMDA~\citep{Thoppilan:22short}, PaLM~\citep{Chowdhery:22short}, Llama~\citep{Touvron:23arxiv}, and Gemini~\citep{GeminiTeam:23arxiv,Reid:24arxiv}.

\section{Experiments}
\label{sec:Exp}

This section presents an empirical assessment of xLSTM in comparison with established approaches, with a particular emphasis on language modeling. Section~\ref{sec:ExpSyntethic} explores the unique strengths of xLSTM on synthetic tasks. 
Section~\ref{sec:ExpComparison} provides a comparative analysis of the validation perplexities achieved by several contemporary language modeling techniques trained on 15B tokens from SlimPajama~\citep{Soboleva:23}, including ablation studies for xLSTM. Furthermore, the scaling properties of all evaluated methods are examined in line with the methodologies described in \citet{Kaplan:20} and \citet{Brown:20short}.
In Section~\ref{sec:ExpLanguage},
a comprehensive language modeling experiment is performed: xLSTM and the highest-performing models from Section~\ref{sec:ExpComparison} are trained using 300B tokens from SlimPajama~\citep{Soboleva:23}. The evaluation comprises (1) measuring the ability of the models to generalize to extended context lengths, (2) analysis of validation perplexity and downstream task performance~\citep{Sutawika:24}, (3) testing on 571 textual domains included in the PALOMA language benchmark~\citep{Magnusson:23arxivshort}, and (4) examining scaling behaviors across methods with a dataset twenty times larger than in the prior experiments.

\label{sec:xlstm_blocks_notation}
Throughout all experimental setups, the notation xLSTM[$a$:$b$] is employed to denote the proportion $a/b$ between mLSTM-based and sLSTM-based xLSTM blocks. For instance, xLSTM[7:1] indicates a configuration with eight blocks, comprising seven mLSTM-based and one sLSTM-based block. In the frequently used setup with 48 total blocks, this ratio yields 42 mLSTM-based and 6 sLSTM-based blocks. Additionally, each experiment incorporates pre up-projection blocks for mLSTM and post up-projection blocks for sLSTM.

\subsection{Synthetic Tasks and Long Range Arena}
\label{sec:ExpSyntethic}
To begin, we evaluate the impact of xLSTM's exponential gating mechanism combined with memory mixing on formal language tasks~\citep{Deletang:23}. Next, we investigate the utility of the xLSTM's matrix memory design using the Multi-Query Associative Recall benchmark~\citep{Arora:23arxiv}. Lastly, we measure xLSTM's ability to handle extended sequence inputs by conducting experiments in the Long Range Arena~\citep{Tay:21}.

\paragraph{Evaluation of xLSTM's Exponential Gating with Memory Mixing.} We assess the effectiveness of xLSTM's exponential gating combined with memory mixing, a mechanism hypothesized to enhance its ability to manage state tracking tasks~\citep{Merrill:24, Merrill:23}. Building upon the formal language benchmarks of \citet{Deletang:23}, we adapt and broaden these tasks to allow for training across varying sequence lengths, thus enabling extrapolation to longer sequences. Appendix~\ref{sec:appExpSynthetic-formalLang} provides comprehensive details regarding task specifications and further experimental outcomes. 

Our analysis benchmarks xLSTM against several alternative approaches, such as Transformers, State Space Models, and traditional Recurrent Neural Networks. We measure performance based on token-level accuracy, focusing exclusively on tokens pertinent to the given tasks. Scores are normalized between 0 (indicating chance performance) and 1 (denoting flawless performance). The models evaluated include various two-block configurations: xLSTM[0:1] (containing solely sLSTM), xLSTM[1:0] (containing only mLSTM), xLSTM[1:1], Llama, Mamba, RWKV, Retention, Hyena, LSTM, as well as LSTM integrated within Transformer blocks (LSTM (Block)). The experimental findings are depicted in Figure~\ref{fig:formal-main}.

Our results indicate that models such as Transformers and State Space Models lacking memory mixing and, hence, state tracking capability, fail to address tasks such as regular grammars exemplified by the parity task. This observation corroborates previous literature indicating that Transformers and State Space models are fundamentally limited in comparison to RNNs with respect to stateful computation~\citep{Merrill:24, Merrill:23, Deletang:23}.

\paragraph{Test of xLSTM's Memory Capacities on Associative Recall Tasks.} This study investigates the memory capacity of xLSTM’s novel matrix-based memory using the Multi-Query Associative Recall task~\citep{Arora:23arxiv}. In this setup, each input sequence comprises key-value pairs, sampled at random from a sizeable vocabulary, that the model is required to store and later recall. We introduce greater complexity to the canonical version by scaling the number of key-value pairs up to 256 and increasing the context window to as many as 2048 tokens, thereby providing a more extensive evaluation of the memory abilities across models. Our assessment includes 2-block variants of Llama, Mamba, RWKV-5, RWKV-6, xLSTM[1:1], and xLSTM[1:0], with performance measured by the accuracy of pair recall. Given that the memory in Transformers (such as Llama) grows exponentially with respect to the representation size~\citep{Ramsauer:21}, these models serve as the benchmark for this task. Figure~\ref{fig:mqar-main} summarizes the outcomes. 
Among non-Transformer models, xLSTM[1:1] yields the highest recall accuracy, even for models of smaller size. Notably, including the sLSTM block does not adversely affect memory storage; on the contrary, it enhances it, which becomes particularly clear in the challenging scenario with 256 key-value associations. Further findings can be found in Appendix~\ref{sec:appExpSynthetic-mqar}, where extrapolation experiments suggest that the superior memory capacity of xLSTM enables successful generalization to context lengths that surpass those encountered during training.

\paragraph{Test of xLSTM's Long Context Capabilities on Long Range Arena.} To evaluate how xLSTM manages extended sequences and substantial context windows, we benchmark several approaches using the Long Range Arena~\citep{Tay:21}. Across all tasks, xLSTM exhibits reliable and robust results, indicating that its design is highly effective in addressing various challenges associated with processing lengthy contexts.

For more details, see Appendix~\ref{sec:appExpSynthetic-lra}.

\subsection{Method Comparison and Ablation Study}
\label{sec:ExpComparison}

In order to investigate the central question of this study—specifically, the performance of our proposed LSTM variants when scaled for language modeling—we conduct experiments in which xLSTMs, Transformers, State Space Models, and additional architectures are trained on a dataset of 15B tokens from SlimPajama using an identical auto-regressive approach. The resulting models are evaluated against the validation set, and further, ablation analyses are carried out for xLSTM configurations.

\paragraph{Benchmarking xLSTM Against Alternative Approaches.} To facilitate a fair comparison, each model is trained on 15B tokens sourced from SlimPajama~\citep{Soboleva:23}. Performance is gauged based on perplexity over the validation dataset. We assess the following architectures: xLSTM (our proposed method), GPT-3 (Transformer)~\citep{Brown:20short}, Llama (Transformer)~\citep{Touvron:23arxiv}, H3 (SSM)~\citep{Fu:23}, Mamba (SSM)~\citep{Gu:24arxiv}, RWKV-4 (RNN)~\citep{Peng:23arxivshort}, RWKV-5 (RNN)~\citep{Peng:24arxivshort}, RWKV-6 (RNN)~\citep{Peng:24arxivshort}, GLA (linear Transformer)~\citep{Yang:23arxiv}, HGRN (RNN)~\citep{Qin:23}, HGRN2 (RNN)~\citep{Qin:24arxiv}, RetNet (linear Transformer)~\citep{Sun:23arxiv}, Hyena (linear Transformer)~\citep{Poli:23}, as well as xLSTM[1:0] and xLSTM[7:1] variants (see Section~\ref{sec:xlstm_blocks_notation}). 
All models are trained in mixed precision; however, RWKV-5, RWKV-6, GLA, and HGRN2 do not use PyTorch’s automated mixed-precision capabilities (refer to Appendix Section~\ref{sec:appGenTrainProc}). 
Methods are organized into three classes: (a) Transformers, (b) State Space Models (SSMs), and (c) Recurrent Neural Networks (RNNs) including linear Transformer variants. Linear Transformers provide an alternative to traditional Transformer attention with linear mechanisms. Each model configuration matches GPT-3 in terms of parameter count (350M), employing an embedding dimension of 1024 and comprising 24 residual blocks. Notably, GPT-3 is unique in utilizing shared weights between token and output embeddings, resulting in a slightly reduced parameter count. 
As shown in Table~\ref{tab:spaj15b_model_results}, xLSTM achieves the lowest validation perplexity, surpassing the other evaluated approaches. Additional specifics can be found in Appendix~\ref{sec:appExpComparison}. Figure~\ref{fig:temp_sclaw15B} illustrates the scaling trends for this setup, supporting the claim that xLSTM retains its advantages as model scale increases.

\paragraph{Ablation Studies.}
As shown in Table~\ref{tab:spaj15b_model_results} and Figure~\ref{fig:temp_sclaw15B}, xLSTM performs remarkably well in language modeling tasks, particularly when trained on 15B tokens sourced from SlimPajama. To isolate the contributions of each modification transitioning from the standard LSTM to xLSTM, we systematically transform a vanilla LSTM model into its xLSTM counterpart. The process begins by embedding the LSTM layers within pre-LayerNorm residual frameworks. The next stage involves supplementing with a post up-projection block. Lastly, exponential gating mechanisms and matrix memory are incorporated.

The results are shown in Table~\ref{tab:ablstudies} (top). 

The observed enhancements in performance are credited to the synergistic effects of exponential gating and the incorporation of matrix memory, as revealed by the ablation studies. Given the critical role of gating in Recurrent Neural Networks (RNNs) and State Space Models, various gating approaches are systematically ablated for comparison. As shown in Table~\ref{tab:ablstudies} (bottom), allowing each gate to be trainable and directly modulated by the input consistently yields further improvements. Further analyses on separate backbone modules are provided in Appendix~\ref{sec:appAblDetails}.

\subsection{xLSTM as Large Language Model}
\label{sec:ExpLanguage}

This work concludes with extensive language modeling experiments, evaluating xLSTM's suitability as a large language model (LLM). To this end, we expand the training dataset and use 300B tokens from SlimPajama, aligning with the token counts employed for Mamba~\citep{Gu:24arxiv} and Griffin~\citep{De:24arxiv}.

Our comparison includes xLSTM, RWKV-4, Llama, and Mamba—each representing their respective model class as described in Section~\ref{sec:ExpComparison}. RWKV-4 is chosen as the RNN exemplar, since suitable precision settings for RWKV-5, RWKV-6, and HGRN2 were determined only after the commencement of training on the 300B token experiments (refer to Appendix~\ref{sec:appGenTrainProc}).

We train multiple model sizes (125M, 350M, 760M, 1.3B), evaluate all variants for their sequence length extrapolation abilities, and assess their validation set performance. Further evaluation covers downstream task proficiency, language modeling across 571 text domains from the PALOMA benchmark, and a detailed examination of their scaling law characteristics.

\paragraph{Sequence Length Extrapolation.}
\label{sec:spaj300B_ctx_extrapolation}
We begin by evaluating the ability of 1.3B-parameter variants of xLSTM, RWKV-4, Llama, and Mamba to extrapolate to longer input sequences. Models are initially trained on a context length of 2048 tokens and subsequently assessed on sequences with lengths up to 16384 tokens. The performance outcomes are visualized in Figure~\ref{fig:spaj300B_seq_extrapolation}. Unlike the alternative architectures, xLSTM continues to achieve low perplexity scores when handling extended context lengths.

\paragraph{Validation Perplexity and Downstream Tasks.}
\label{sec:spaj_downstream_eval}
Next, across all scales, we assess xLSTM, RWKV-4, Llama, and Mamba on their ability to predict the next token using the SlimPajama validation dataset, as well as their effectiveness on downstream benchmarks designed to evaluate common sense reasoning capabilities.

The third column in Table~\ref{tab:lmeval} presents the validation set perplexity achieved by various methods. For all model sizes, xLSTM[1:0] and xLSTM[7:1] consistently yield the lowest perplexity on the validation set. The remaining columns in Table~\ref{tab:lmeval} show results on a range of downstream tasks. Across almost all tasks and model sizes, xLSTM demonstrates superior performance compared to the alternatives; Mamba outperforms xLSTM only in select instances on the ARC task. Further information can be found in Appendix~\ref{sec:appExpLanguage}.

\paragraph{Performance on PALOMA Language Tasks.} In addition, we evaluate the capabilities of xLSTM, RWKV-4, Llama, and Mamba across various model configurations in the context of PALOMA language tasks~\citep{Magnusson:23arxivshort}. Specifically, next token prediction is assessed by calculating perplexity across 571 distinct text domains, encompassing content from sources like nytimes.com to the Reddit subreddit r/depression. 
Table~\ref{tab:ppleval} displays perplexity values, split into language modeling (shown in the first seven columns) and detailed domain benchmarks (presented in the final five columns). xLSTM[1:0] demonstrates superior performance compared to xLSTM[7:1] on these tasks.

In this evaluation, xLSTM[1:0] outperforms Mamba in 568 out of 571 domains (99.5\%), Llama in 486 domains (85.1\%), and RWKV-4 in 570 domains (99.8\%) in terms of lower perplexity scores. Further information can be found in Appendix~\ref{sec:appExpLanguage}.

\paragraph{Scaling Laws.} In the fourth part of our analysis, we investigate the power-law scaling dynamics, which provide a basis for projecting model performance as the size increases~\citep{Kaplan:20,Brown:20short}. The scaling relationships are illustrated in Figure~\ref{fig:temp_sclaw300B}. Although each model exhibits comparable scaling trends, they differ in their performance baselines. RWKV-4 exhibits the lowest performance, with Llama and Mamba following in ascending order. xLSTM outperforms Mamba, maintaining a performance gap similar to that observed between Mamba and Llama. These scaling results suggest that, as model size grows, xLSTM is expected to continue outperforming Transformers and State-Space models.

\textbf{Generation Times and Maximal Throughput.} In our final set of evaluations, depicted in Figure~\ref{fig:inference_time_generation} and Figure~\ref{fig:inference_time_decoding_throughput} (left), we investigate the text generation latency and peak throughput of several xLSTM variants at the 1.3B parameter scale. Our analysis includes comparisons to equivalently-sized implementations of Mamba, Llama, and RWKV sourced from HuggingFace, with the Llama benchmark utilizing a static key-value cache. It should be noted that, during these experiments, it was not feasible to compile the full cache for the Transformer Model, nor was compilation of the Mamba model using \texttt{torch.compile} successful. In the generation benchmarks, all models are assessed with a batch size of 1 and a pre-fill of 16 tokens, which particularly advantages the Transformer-based Llama.

Figure~\ref{fig:inference_time_generation} demonstrates that xLSTM, along with other recurrent architectures such as Mamba and RWKV-4, exhibits linear growth in generation time, in stark contrast with the quadratic time complexity shown by Llama. Throughput analyses were also conducted using varying batch sizes and pre-fill parameters for Llama. The results in Figure~\ref{fig:inference_time_decoding_throughput} (right) indicate that xLSTM leverages its fixed memory requirements to efficiently accommodate substantially larger batch sizes than Llama, thereby attaining superior maximal throughput.

\section{Limitations}

(i) Unlike mLSTM, sLSTM’s memory mixing mechanism restricts the possibility of parallelization, thereby preventing efficient parallel processing. Nonetheless, we have created an optimized CUDA kernel for sLSTM, which achieves a runtime less than twice that of our parallel mLSTM implementation. (ii) The current CUDA kernels for mLSTM lack optimization, causing their execution speed to be roughly four times slower than FlashAttention or the scan methodology employed in Mamba. Enhanced CUDA kernels similar to those of FlashAttention could substantially improve performance. (iii) mLSTM’s matrix memory introduces considerable computational demands, as it involves the manipulation of $d \times d$ matrices. However, both memory updates and retrievals are parameter-free and amenable to standard parallel matrix operations, meaning the impact on runtime remains limited. (iv) Careful selection of initial values for the forget gates is critical. (v) Because the matrix memory does not depend on sequence length, extending the sequence size may result in memory saturation for longer contexts. Nevertheless, contexts up to 16k tokens do not seem to pose a problem, as discussed in Section~\ref{sec:spaj300B_ctx_extrapolation}. (vi) Given the substantial computational resources required for large-scale language modeling, our work did not include thorough architectural or hyperparameter optimization, particularly in the context of more extensive xLSTM variants. We believe that comprehensive tuning will be essential to fully leverage the capabilities of xLSTM.

\section{Conclusion}
Our investigation has provided a partial answer to the central question: What progress can be achieved in language modeling by scaling LSTM architectures to billions of parameters? To date, the evidence suggests: ``At a minimum, results are comparable to established approaches such as Transformers and State Space Models.'' We introduced xLSTM, augmenting LSTM with exponential gating mechanisms incorporating memory mixing and an updated memory architecture. The empirical results demonstrate that xLSTM models yield strong performance on language modeling benchmarks, outperforming or rivaling the leading methods including Transformers and State Space Models. Observations of scaling behavior imply that further increasing xLSTM capacity will position them as formidable contenders against the latest Large Language Models predominantly employing Transformer frameworks. Additionally, xLSTM’s architecture shows promise for substantial advancements in related domains such as Reinforcement Learning, Time Series Forecasting, and the computational modeling of physical phenomena.

\end{document}